\textbf{Phần 1. Trắc nghiệm lựa chọn}
\setcounterpageref{ex}{0}

\begin{ex} %[2D1V06]
Một bác thợ mộc cưa một khúc gỗ dài $12\text{ m}$ thành các đoạn nhỏ dài $3\text{ m}$. Hỏi bác thợ phải thực hiện bao nhiêu nhát cưa?
\choice
{$4$ nhát cưa}
{\True $3$ nhát cưa}
{$5$ nhát cưa}
{$2$ nhát cưa}
\loigiai{
Số đoạn gỗ cưa được là: $12 : 3 = 4$ (đoạn).\nSố nhát cưa cần thực hiện là: $4 - 1 = 3$ (nhát cưa).
\textbf{Đáp án của câu hỏi là B}
}
\end{ex}

\begin{ex} %[2D1V06]
Người ta trồng cây dọc theo một lối đi dài $20\text{ m}$, hai đầu đường đều có cây. Biết khoảng cách giữa hai cây liên tiếp là $5\text{ m}$. Hỏi có tất cả bao nhiêu cây?
\choice
{\True $5$ cây}
{$4$ cây}
{$6$ cây}
{$7$ cây}
\loigiai{
Số khoảng cách là: $20 : 5 = 4$ (khoảng cách).\nVì hai đầu đường đều trồng cây nên số cây là: $4 + 1 = 5$ (cây).
\textbf{Đáp án của câu hỏi là A}
}
\end{ex}

\begin{ex} %[2D1V06]
Tôi là một số có hai chữ số. Nếu lấy tôi trừ đi $15$ rồi cộng thêm $24$ thì được số tròn chục lớn nhất có hai chữ số. Tôi là số nào?
\choice
{$79$}
{$90$}
{\True $81$}
{$85$}
\loigiai{
Số tròn chục lớn nhất có hai chữ số là 90.\nSố cần tìm là: $90 - 24 + 15 = 66 + 15 = 81$.
\textbf{Đáp án của câu hỏi là C}
}
\end{ex}

\begin{ex} %[2D1V06]
Hình vẽ gồm một đoạn thẳng dài $35\text{ cm}$ bị cắt bớt một đoạn dài $1\text{ dm } 2\text{ cm}$. Đoạn thẳng còn lại dài bao nhiêu?
\choice
{$25\text{ cm}$}
{$21\text{ cm}$}
{$22\text{ cm}$}
{\True $23\text{ cm}$}
\loigiai{
Đổi $1\text{ dm } 2\text{ cm} = 12\text{ cm}$.\nĐoạn thẳng còn lại dài: $35 - 12 = 23\text{ (cm)}$.
\textbf{Đáp án của câu hỏi là D}
}
\end{ex}

\begin{ex} %[2D1V06]
Tìm số tự nhiên $y$ biết: $y + 16 < 20$. Các giá trị thích hợp của $y$ là
\choice
{$1, 2, 3, 4$}
{\True $0, 1, 2, 3$}
{$0, 1, 2$}
{$1, 2, 3$}
\loigiai{
Ta có: $y < 20 - 16 \Rightarrow y < 4$. Vì $y$ là số tự nhiên nên $y \in \{0, 1, 2, 3\}$.
\textbf{Đáp án của câu hỏi là B}
}
\end{ex}

\begin{ex} %[2D1V06]
Một giá sách có hai ngăn đựng tất cả $78$ quyển sách. Nếu chuyển $8$ quyển sách từ ngăn trên xuống ngăn dưới thì hai ngăn có số sách bằng nhau. Số sách ban đầu của ngăn trên là
\choice
{\True $47$ quyển}
{$45$ quyển}
{$39$ quyển}
{$49$ quyển}
\loigiai{
Khi hai ngăn bằng nhau, mỗi ngăn có: $78 : 2 = 39$ (quyển).\nSố sách ban đầu của ngăn trên là: $39 + 8 = 47$ (quyển).
\textbf{Đáp án của câu hỏi là A}
}
\end{ex}

\begin{ex} %[2D1V06]
Đồng hồ chỉ $7$ giờ sáng. Kim giờ và kim phút tạo thành một góc. Sau $30$ phút nữa thì đồng hồ chỉ mấy giờ?
\choice
{$8$ giờ}
{$7$ giờ $15$ phút}
{\True $7$ giờ $30$ phút}
{$6$ giờ $30$ phút}
\loigiai{
Sau 30 phút, thời gian là $7$ giờ $+ 30$ phút $= 7$ giờ $30$ phút.
\textbf{Đáp án của câu hỏi là C}
}
\end{ex}

\begin{ex} %[2D1V06]
Cho ba số có tổng bằng $95$. Biết số thứ nhất là $23$, số thứ hai hơn số thứ nhất $8$ đơn vị. Số thứ ba là
\choice
{$39$}
{$43$}
{$45$}
{\True $41$}
\loigiai{
Số thứ hai là: $23 + 8 = 31$.\nSố thứ ba là: $95 - (23 + 31) = 95 - 54 = 41$.
\textbf{Đáp án của câu hỏi là D}
}
\end{ex}

\begin{ex} %[2D1V06]
Tìm số lớn nhất có hai chữ số mà chữ số hàng chục hơn chữ số hàng đơn vị $4$ đơn vị.
\choice
{$84$}
{\True $95$}
{$73$}
{$94$}
\loigiai{
Để số lớn nhất, chọn hàng chục là $9$. Chữ số hàng đơn vị là $9 - 4 = 5$. Vậy số đó là $95$.
\textbf{Đáp án của câu hỏi là B}
}
\end{ex}

\begin{ex} %[2D1V06]
Trong một phép trừ, nếu bớt ở số bị trừ đi $7$ đơn vị và giữ nguyên số trừ thì hiệu sẽ
\choice
{\True Giảm đi $7$ đơn vị}
{Tăng thêm $7$ đơn vị}
{Không đổi}
{Giảm đi $14$ đơn vị}
\loigiai{
Khi bớt ở số bị trừ đi bao nhiêu đơn vị thì hiệu cũng giảm đi bấy nhiêu đơn vị.
\textbf{Đáp án của câu hỏi là A}
}
\end{ex}

\begin{ex} %[2D1V06]
Cho hình chữ nhật có chiều dài $25\text{ cm}$, chiều rộng ngắn hơn chiều dài $11\text{ cm}$. Chiều rộng của hình chữ nhật đó là
\choice
{$15\text{ cm}$}
{$16\text{ cm}$}
{\True $14\text{ cm}$}
{$12\text{ cm}$}
\loigiai{
Chiều rộng hình chữ nhật là: $25 - 11 = 14\text{ (cm)}$.
\textbf{Đáp án của câu hỏi là C}
}
\end{ex}

\begin{ex} %[2D1V06]
Bạn Bình đếm các số tròn chục từ $10$ đến $90$. Hỏi bạn Bình đã đếm được bao nhiêu chữ số $0$?
\choice
{$8$ chữ số $0$}
{$10$ chữ số $0$}
{$7$ chữ số $0$}
{\True $9$ chữ số $0$}
\loigiai{
Các số là: 10, 20, 30, 40, 50, 60, 70, 80, 90. Mỗi số có đúng 1 chữ số 0 nên có 9 chữ số 0.
\textbf{Đáp án của câu hỏi là D}
}
\end{ex}

\textbf{Phần 2. Trắc nghiệm đúng sai}
\setcounterpageref{ex}{0}

\begin{ex} %[2D1V06]
Xét tính đúng sai của các bài toán tư duy tổng hợp:
\choiceTF
{\True Một sợi dây dài 10 m bị cắt thành 5 đoạn bằng nhau thì cần 4 nhát cắt.}
{Nếu hôm nay là thứ Hai ngày 4 thì thứ Hai tuần sau là ngày 10.}
{\True Số nhỏ nhất có hai chữ số mà tổng hai chữ số bằng 9 là số 18.}
{\True Tổng của số liền trước và số liền sau của 50 là 100.}
\loigiai{
\begin{itemize}
    \item Ý a) Đúng: Số nhát cắt = Số đoạn - 1 = $5 - 1 = 4$ nhát.
    \item Ý b) Sai: Một tuần có 7 ngày, thứ Hai tuần sau là ngày $4 + 7 = 11$.
    \item Ý c) Đúng: Chọn chữ số hàng chục nhỏ nhất là 1, hàng đơn vị là $9 - 1 = 8$. Số đó là 18.
    \item Ý d) Đúng: $49 + 51 = 100$. Đúng.
\end{itemize}
\begin{itemize}
    \item a) \textbf{Đúng}.
    \item b) \textbf{Sai}.
    \item c) \textbf{Đúng}.
    \item d) \textbf{Đúng}.
\end{itemize}
}
\end{ex}

\begin{ex} %[2D1V06]
Cho dãy số: 2, 5, 8, 11, 14, ...:
\choiceTF
{\True Quy luật của dãy số là mỗi số đứng sau bằng số đứng trước cộng thêm 3.}
{Số hạng thứ 7 của dãy số là số 19.}
{\True Các số trong dãy số trên khi chia cho 3 đều dư 2.}
{Số 30 là một số thuộc dãy số trên.}
\loigiai{
\begin{itemize}
    \item Ý a) Đúng: $5 - 2 = 3$, $8 - 5 = 3$.
    \item Ý b) Sai: Số thứ 6 là 17, số thứ 7 là $17 + 3 = 20$ (không phải 19).
    \item Ý c) Đúng: $2 : 3$ dư 2, $5 : 3$ dư 2, $8 : 3$ dư 2,... Khẳng định đúng.
    \item Ý d) Sai: Các số trong dãy chia 3 dư 2, còn 30 chia hết cho 3 nên 30 không thuộc dãy.
\end{itemize}
\begin{itemize}
    \item a) \textbf{Đúng}.
    \item b) \textbf{Sai}.
    \item c) \textbf{Đúng}.
    \item d) \textbf{Sai}.
\end{itemize}
}
\end{ex}

\textbf{Phần 3. Trả lời ngắn}
\setcounterpageref{ex}{0}

\begin{ex} %[2D1V06]
Một bác thợ cưa một cây tre thành 6 đoạn. Mỗi nhát cưa mất 2 phút. Hỏi bác thợ cưa xong cây tre trong bao nhiêu phút?
\shortans{10}
\loigiai{
\textbf{Đáp số:} $10$\\
Để có 6 đoạn cần $6 - 1 = 5$ nhát cưa. Thời gian cưa là: $5 \times 2 = 10$ phút.
}
\end{ex}

\begin{ex} %[2D1V06]
Tìm số bé nhất có hai chữ số mà chữ số hàng đơn vị hơn chữ số hàng chục 6 đơn vị.
\shortans{17}
\loigiai{
\textbf{Đáp số:} $17$\\
Chọn chữ số hàng chục nhỏ nhất là 1, chữ số hàng đơn vị là $1 + 6 = 7$. Số đó là 17.
}
\end{ex}

\begin{ex} %[2D1V06]
An và Bình có tất cả 46 viên bi. An cho Bình 5 viên bi thì hai bạn có số bi bằng nhau. Hỏi lúc đầu An có bao nhiêu viên bi?
\shortans{28}
\loigiai{
\textbf{Đáp số:} $28$\\
Khi bằng nhau mỗi bạn có $46 : 2 = 23$ viên. Lúc đầu An có: $23 + 5 = 28$ viên.
}
\end{ex}

\begin{ex} %[2D1V06]
Tìm một số biết rằng lấy số lớn nhất có hai chữ số trừ đi số đó rồi trừ tiếp 25 thì được 30.
\shortans{44}
\loigiai{
\textbf{Đáp số:} $44$\\
Số đó là: $99 - 25 - 30 = 44$.
}
\end{ex}

\begin{ex} %[2D1V06]
Lớp 2A xếp thành 4 hàng, mỗi hàng có 8 học sinh và thừa ra 3 bạn. Hỏi lớp 2A có tất cả bao nhiêu học sinh?
\shortans{35}
\loigiai{
\textbf{Đáp số:} $35$\\
$4 \times 8 + 3 = 32 + 3 = 35$ học sinh.
}
\end{ex}

\begin{ex} %[2D1V06]
Tìm x biết: $x + 14 = 85 - 21$.
\shortans{50}
\loigiai{
\textbf{Đáp số:} $50$\\
Vế phải: $85 - 21 = 64$. Do đó $x = 64 - 14 = 50$.
}
\end{ex}
